\begin{proposition}[Euler number with a higher-genus ruled resolution]\hypertarget{tag-02YJ}{}\label{night:P4R.7}\label{stmt:P4R.7}
\StatementRecord{night:P4R.7}{proved}{02YJ}{\SeriesRef{night:P4.13}, \SeriesRef{night:P4R.1}, \SeriesRef{night:P4R.3}, \SeriesRef{night:P4R.6}}

Let $f:Y\to B$ satisfy the torus-fibration hypotheses of
\cref{night:P4.12}, and let $mS$ be a fibre with irreducible,
reduced, singular support. Let $\nu:N\to S$ be the normalization
and $\pi:T\to N$ its minimal resolution. Suppose that $T$ is
birationally ruled over a curve of genus at least two. Then
$e(S)\ge0$, including when $N$ is singular.
\end{proposition}
\begin{proof}
\begin{enumerate}[label=\arabic*.,leftmargin=*]
\item Put $\rho=\nu\circ\pi$, and take the effective integral
divisor $D$ from \cref{night:P4R.6}. Let $\beta:T\to T_0$ be
a sequence of point blowdowns to a minimal ruled surface and
put $D_0=\beta_*D$. The canonical divisor identity makes
$K_T+D$ torsion, so $K_{T_0}+D_0\equiv0$. The intersection
of $D_0$ with a ruling fibre is two. The adjunction and
Riemann--Hurwitz calculation in the proof of
\cref{night:P4R.3} uses these numerical and effectiveness
hypotheses and gives
\[
 D_0=2H+\sum_j a_jF_j,
 \qquad a_j\in\mathbb Z_{>0},
\]
where $H$ is a smooth section and the $F_j$ are distinct
ruling fibres.

At a reverse point blowup, a centre of multiplicity $\mu$
gives exceptional coefficient $\mu-1$. At a crossing,
$\mu\ge2$, so the exceptional curve subdivides the edge of
the reduced divisor. At a smooth point, it adds a rational
leaf when $\mu\ge2$ and leaves the reduced divisor unchanged
when $\mu=1$. Thus $A=D_{\rm red}$ is the strict transform
of $H$ with rational trees attached, and is connected.
Collapsing those trees to their attachment points induces
an isomorphism on fundamental groups by van Kampen.
The ruling and point blowdowns also induce isomorphisms
on fundamental groups. Their composite restricts on the
strict transform of $H$ to the section isomorphism with
the ruling base. Therefore the inclusion
$j:A\hookrightarrow T$ induces an isomorphism
$j_*:\pi_1(A)\to\pi_1(T)$.

\item Put $C=\operatorname{Sing}(S)_{\rm red}$.
Proposition~\ref{night:P4R.6} gives $A=\rho^{-1}(C)_{\rm red}$,
and $\rho$ identifies $T\setminus A$ with $S_{\rm reg}$.
In the quotient $T\cup_A C$, every fibre of $\rho$ over
$c\in C$ is identified with $c$. The induced map to $S$
is consequently a continuous bijection from a compact
space to a Hausdorff space, hence a homeomorphism.
This quotient uses the whole proper map $\rho|_A:A\to C$,
including the exceptional components contracted to points.

A triangulation of the analytic pair $(T,A)$ makes
$A\subset T$ a subcomplex and a cofibration
\cite[Theorem~2, p.~464]{Lojasiewicz1964}.
The associated double mapping cylinder has the homotopy
type of this pushout. Van Kampen therefore gives
\[
 \pi_1(S)\simeq
 \pi_1(T)*_{\pi_1(A)}\pi_1(C)\simeq\pi_1(C).
\]
For the second isomorphism, the inverse map is identity
on $\pi_1(C)$ and is $(\rho|_A)_*j_*^{-1}$ on
$\pi_1(T)$; these maps agree on the amalgamating group.
Proposition~\ref{night:P4R.1} makes $\pi_1(C)$ virtually abelian.

Let $c$ be the number of irreducible
components of $C$, let $g_i$ be their normalization genera,
and let $b$ be the cycle rank of its normalization incidence
graph. The calculation in \cref{night:P4R.3} gives
\[
 \pi_1(C)\simeq
   \Bigl(*_i\pi_1(C_i^{\rm norm})\Bigr)*F_b,
 \qquad e(C)=c+1-2\sum_i g_i-b.
\]
Virtual abelianness permits either all $g_i=0$ with
$b\le1$, or one $g_i=1$ with $b=0$. These give respectively
$e(C)=c+(1-b)\ge0$ and $e(C)=c-1\ge0$.
Finally \cref{night:P4.13}, applied to this same fibre $mS$,
gives $e(S_{\rm reg})\ge0$. Additivity yields
\[
 e(S)=e(S_{\rm reg})+e(C)\ge0.
\]
\end{enumerate}
\end{proof}